% ===========================================================================
\section{Conjunto de instrucciones}
\label{sec:isa}

El procesador implementa el subset de RV32I que pide la consigna más una instrucción de parada
(\repo{docs/spec/isa.md}):

\begin{center}
\small
\begin{tabular}{@{}l l l@{}}
  \toprule
  \textbf{Tipo} & \textbf{Formato} & \textbf{Instrucciones} \\
  \midrule
  Aritméticas y lógicas    & R & \texttt{add sub sll srl sra and or xor slt sltu} \\
  Con inmediato            & I & \texttt{addi andi ori xori slti sltiu slli srli srai} \\
  Cargas                   & I & \texttt{lb lh lw lbu lhu} \\
  Almacenamientos          & S & \texttt{sb sh sw} \\
  Saltos condicionales     & B & \texttt{beq bne} \\
  Saltos incondicionales   & J, I & \texttt{jal jalr} \\
  Inmediato superior       & U & \texttt{lui} \\
  Parada                   & I & \texttt{halt} (= \texttt{ebreak}, \texttt{0x00100073}) \\
  \bottomrule
\end{tabular}
\end{center}

\diagrama{formatos.png}{Formatos de instrucción}{Los seis formatos de RISC-V y, para cada uno con
inmediato, de qué bits de la instrucción sale cada bit del inmediato de 32 bits. Los colores
identifican el bit de origen y son los mismos en todos los formatos.}{fig:formatos}

\subsection{La instrucción de parada}
\label{sec:halt}

La consigna pide una instrucción que detenga el procesador con el pipeline vacío. Se codificó como
\texttt{ebreak} (\texttt{0x00100073}) porque no choca con ninguna instrucción implementada, el
ensamblador estándar la genera y no exige herramientas propias (\decision{001}).

Viaja por el pipeline como una instrucción sin efectos que solo lleva su bit \texttt{halt}.
Mientras haya un HALT en cualquier etapa, IF deja de buscar:

\begin{lstlisting}[style=verilog,numbers=none]
stop_fetch  = id_is_halt | id_ex_halt | ex_mem_halt | mem_wb_halt | halted;
flush_if_id = redirect | stop_fetch;
\end{lstlisting}

\texttt{stop\_fetch} congela el PC y además convierte en burbuja lo que entra a IF/ID. Hacen falta
las dos cosas: cuando el HALT está en ID, la BRAM ya está leyendo la dirección siguiente y esa
palabra aparecería en IF/ID en el flanco siguiente. La parada mira las cuatro etapas porque el HALT
está en ID un solo ciclo. Cuando llega a WB se levanta \texttt{halted} y el pipeline queda vacío
(los cuatro \texttt{valid} en 0).

Hay dos casos borde que la verificación cubre con programas propios:
\begin{itemize}[nosep]
  \item \textbf{HALT en el camino equivocado}: si un salto anterior se toma, el \emph{flush}
        convierte al HALT en burbuja, su bit desaparece y \texttt{stop\_fetch} se apaga solo
        (\repo{sw/halt_camino_equivocado.s}).
  \item \textbf{HALT retenido por un stall}: \texttt{0x00100073} tiene un 1 en el campo
        \texttt{rs2}, así que un \texttt{lw x1} justo antes dispara la detección de carga-uso.
        Como el \emph{enable} tiene prioridad sobre el \emph{flush}, IF/ID conserva el HALT
        (\repo{sw/halt_stall.s}).
\end{itemize}

\subsection{Instrucciones fuera del subset}

La unidad de control decodifica por \emph{opcode} y compara la palabra completa solo para el HALT
(\decision{024}). Un opcode desconocido se comporta como NOP. El ensamblador y el debugger
rechazan cualquier palabra que no esté en \repo{isa.md} antes de cargarla (\decision{025}), así
que esos casos no llegan a la placa desde las herramientas del proyecto.
